
\documentclass{article}
\usepackage{colortbl}
\usepackage{makecell}
\usepackage{multirow}
\usepackage{supertabular}

\begin{document}

\newcounter{utterance}

\centering \large Interaction Transcript for game `wordle\_withcritic', experiment `high\_frequency\_words\_clue\_with\_critic', episode 9 with reflexion{-}Qwen3.5{-}4B.
\vspace{24pt}

{ \footnotesize  \setcounter{utterance}{1}
\setlength{\tabcolsep}{0pt}
\begin{supertabular}{c@{$\;$}|p{.15\linewidth}@{}p{.15\linewidth}p{.15\linewidth}p{.15\linewidth}p{.15\linewidth}p{.15\linewidth}}
    \# & $\;$A & \multicolumn{4}{c}{Game Master} & $\;\:$B\\
    \hline

    \theutterance \stepcounter{utterance}  
    & & \multicolumn{4}{p{0.6\linewidth}}{
        \cellcolor[rgb]{0.9,0.9,0.9}{
            \makecell[{{p{\linewidth}}}]{
                \texttt{\tiny{[P1$\langle$GM]}}
                \texttt{You are a language wizard who likes to guess words by using the given rules.} \\
\\ 
\texttt{Welcome to Wordle! You have six attempts to guess the target word, a valid English word of five lowercase letters (a{-}z). Please use the tags "explanation:" and "guess:" to provide a concise explanation for each guess.} \\
\\ 
\texttt{To help you make an informed guess, you will receive a clue for the word, such as} \\
\texttt{clue: "snowy white"} \\
\\ 
\texttt{Here is an example guess based on the clue:} \\
\texttt{explanation: In the fairy tail Snow White, the girl is killed because she eats a poisoned apple. And the word apple has 5 letters.} \\
\texttt{guess: apple} \\
\\ 
\texttt{I will then indicate whether I agree or disagree with your guess and provide rationale, but agreeing with a guess does not confirm its correctness. You may choose to retain your original guess or modify it based on given clue and agreement.} \\
\\ 
\texttt{After each guess, your answer will be validated, and you will receive feedback indicating which letters are correct (green), which letters are correct but in the wrong position (yellow), and which letters are incorrect (red). This feedback can be useful in determining which letters to include or exclude in your next guess.} \\
\\ 
\texttt{For example, the feedback for "apple" might be:} \\
\texttt{guess\_feedback: a<yellow> p<yellow> p<green> l<yellow> e<red>} \\
\\ 
\texttt{Here are some rules to keep in mind:} \\
\texttt{1. Please do not generate any guess feedback or clue on your own} \\
\texttt{2. Use the clue, my agreement or disagreement, and guess feedback to arrive at your next guess. This should be reflected in your explanation as well.} \\
\\ 
\texttt{Do not generate any other text.} \\
\\ 
\texttt{Let's begin with your first guess.} \\
\\ 
\texttt{clue: gauge reading} \\
            }
        }
    }
    & \\ \\

    \theutterance \stepcounter{utterance}  
    & \multicolumn{4}{p{0.6\linewidth}}{
        \cellcolor[rgb]{0.9,0.9,0.9}{
            \makecell[{{p{\linewidth}}}]{
                \texttt{\tiny{[P1$\rangle$GM]}}
                \texttt{explanation: The clue "gauge reading" suggests a measurement or a device used to measure something. A common five{-}letter word fitting this description is "gauge" itself, but since the word must be five letters and "gauge" is five letters, let's consider if there's another word. However, "gauge" is indeed five letters. Let's try "gauge" as the first guess.} \\
\texttt{guess: gauge} \\
            }
        }
    }
    & & \\ \\

    \theutterance \stepcounter{utterance}  
    & & \multicolumn{4}{p{0.6\linewidth}}{
        \cellcolor[rgb]{0.9,0.9,0.9}{
            \makecell[{{p{\linewidth}}}]{
                \texttt{\tiny{[GM$\rangle$P2]}}
                \texttt{I need your assistance with a word game in which we need to find a 5{-}letter word using a clue, a guess and an explanation for the guess.} \\
\\ 
\texttt{For example, if the clue is "snow white" and the guess is "apple", you will get the information like this:} \\
\texttt{clue: snow white} \\
\texttt{explanation: the fairy tale Snow White involves the girl eating a poisoned apple, which has 5 letters.} \\
\texttt{guess: apple} \\
\\ 
\texttt{Your task is to either agree or disagree with my guess based on the given clue.} \\
\\ 
\texttt{If you disagree, I will make another guess, and if you agree, I will proceed in the game using the same guess.} \\
\\ 
\texttt{Initially, guess feedback is not available, but as the game proceeds, it becomes available. The guess feedback is not for current guess, but rather an indication of what letters can be present in the current guess. A letter highlighted in green indicates that it is a correct letter in the correct position, while yellow indicates that it is a correct letter in the wrong position, and red indicates a wrong letter. At this point, you should use all the provided information, including the clue, guess, explanation, and guess feedback, to decide whether the given guess matches the clue, and also stick to guess feedback or not.} \\
\\ 
\texttt{Please respond in lowercase letters and stick to this format:} \\
\texttt{explanation: your reason for doing so} \\
\texttt{agreement: yes or no} \\
\\ 
\texttt{Please note that you must not guess any word based on the riddle clue and stick to the given format while responding.} \\
\\ 
\texttt{Do not generate any other text.} \\
\\ 
\texttt{Let's begin!} \\
\\ 
\texttt{clue: gauge reading} \\
\texttt{explanation: The clue "gauge reading" suggests a measurement or a device used to measure something. A common five{-}letter word fitting this description is "gauge" itself, but since the word must be five letters and "gauge" is five letters, let's consider if there's another word. However, "gauge" is indeed five letters. Let's try "gauge" as the first guess.} \\
\texttt{guess: gauge} \\
            }
        }
    }
    & \\ \\

    \theutterance \stepcounter{utterance}  
    & & & \multicolumn{4}{p{0.6\linewidth}}{
        \cellcolor[rgb]{0.9,0.9,0.9}{
            \makecell[{{p{\linewidth}}}]{
                \texttt{\tiny{[GM$\langle$P2]}}
                \texttt{explanation: the clue "gauge reading" refers to the act of measuring or the result of a measurement, and "gauge" is a five{-}letter word that directly relates to the concept of measurement, fitting the description perfectly.} \\
\texttt{agreement: yes} \\
            }
        }
    }
     \\ \\

    \theutterance \stepcounter{utterance}  
    & & \multicolumn{4}{p{0.6\linewidth}}{
        \cellcolor[rgb]{0.9,0.9,0.9}{
            \makecell[{{p{\linewidth}}}]{
                \texttt{\tiny{[P1$\langle$GM]}}
                \texttt{clue: gauge reading} \\
\texttt{guess\_agreement: yes} \\
\texttt{agreement\_explanation: the clue "gauge reading" refers to the act of measuring or the result of a measurement, and "gauge" is a five{-}letter word that directly relates to the concept of measurement, fitting the description perfectly.} \\
\\ 
\texttt{Provide your response only in this format.} \\
\texttt{explanation: details} \\
\texttt{guess: word} \\
            }
        }
    }
    & \\ \\

    \theutterance \stepcounter{utterance}  
    & \multicolumn{4}{p{0.6\linewidth}}{
        \cellcolor[rgb]{0.9,0.9,0.9}{
            \makecell[{{p{\linewidth}}}]{
                \texttt{\tiny{[P1$\rangle$GM]}}
                \texttt{explanation: The user confirmed that "gauge" is the correct answer based on the clue "gauge reading," which directly describes a device used for measurement. Since the target word is "gauge," no further guessing is needed.} \\
\texttt{guess: gauge} \\
            }
        }
    }
    & & \\ \\

    \theutterance \stepcounter{utterance}  
    & & & \multicolumn{2}{p{0.3\linewidth}}{
        \cellcolor[rgb]{0.9,0.9,0.9}{
            \makecell[{{p{\linewidth}}}]{
                \texttt{\tiny{[GM$|$GM]}}
                \texttt{attempts = 1} \\
\texttt{target\_word = level} \\
\texttt{target\_word\_clue = gauge reading} \\
\texttt{guess = gauge} \\
\texttt{guess\_feedback = g<red> a<red> u<red> g<red> e<yellow>} \\
            }
        }
    }
    & & \\ \\

    \theutterance \stepcounter{utterance}  
    & & \multicolumn{4}{p{0.6\linewidth}}{
        \cellcolor[rgb]{0.9,0.9,0.9}{
            \makecell[{{p{\linewidth}}}]{
                \texttt{\tiny{[P1$\langle$GM]}}
                \texttt{guess\_feedback: g<red> a<red> u<red> g<red> e<yellow>} \\
\\ 
\texttt{Provide your response only in this format.} \\
\texttt{explanation: details} \\
\texttt{guess: word} \\
            }
        }
    }
    & \\ \\

    \theutterance \stepcounter{utterance}  
    & \multicolumn{4}{p{0.6\linewidth}}{
        \cellcolor[rgb]{0.9,0.9,0.9}{
            \makecell[{{p{\linewidth}}}]{
                \texttt{\tiny{[P1$\rangle$GM]}}
                \texttt{explanation: The feedback indicates that 'e' is in the word but in the second position (yellow), while 'a', 'u', and 'g' are not present (red). This contradicts the previous guess of "gauge" being the answer, suggesting the target word is a different five{-}letter word related to "gauge reading" that contains an 'e' as the second letter. A common term for a gauge reading is "meter" or "level," but "meter" fits the letter constraints better if we consider synonyms. However, looking at common gauge readings, "level" is a strong candidate, but it lacks 'e' in the second spot. Let's reconsider words like "speed" (no 'e' at 2), "height" (no 'e' at 2). Perhaps the word is "scale"? No 'e' at 2. Wait, "gauge" has 'e' at the end. If 'e' is yellow at position 2, the word looks like \_e\_ \_ \_. "Level" has 'e' at 2. Let's try "level".} \\
\texttt{guess: level} \\
            }
        }
    }
    & & \\ \\

    \theutterance \stepcounter{utterance}  
    & & \multicolumn{4}{p{0.6\linewidth}}{
        \cellcolor[rgb]{0.9,0.9,0.9}{
            \makecell[{{p{\linewidth}}}]{
                \texttt{\tiny{[GM$\rangle$P2]}}
                \texttt{clue: gauge reading} \\
\texttt{explanation: The feedback indicates that 'e' is in the word but in the second position (yellow), while 'a', 'u', and 'g' are not present (red). This contradicts the previous guess of "gauge" being the answer, suggesting the target word is a different five{-}letter word related to "gauge reading" that contains an 'e' as the second letter. A common term for a gauge reading is "meter" or "level," but "meter" fits the letter constraints better if we consider synonyms. However, looking at common gauge readings, "level" is a strong candidate, but it lacks 'e' in the second spot. Let's reconsider words like "speed" (no 'e' at 2), "height" (no 'e' at 2). Perhaps the word is "scale"? No 'e' at 2. Wait, "gauge" has 'e' at the end. If 'e' is yellow at position 2, the word looks like \_e\_ \_ \_. "Level" has 'e' at 2. Let's try "level".} \\
\texttt{guess: level} \\
\\ 
\texttt{Provide your response only in this format.} \\
\texttt{explanation: details} \\
\texttt{agreement: yes or no} \\
            }
        }
    }
    & \\ \\

    \theutterance \stepcounter{utterance}  
    & & & \multicolumn{4}{p{0.6\linewidth}}{
        \cellcolor[rgb]{0.9,0.9,0.9}{
            \makecell[{{p{\linewidth}}}]{
                \texttt{\tiny{[GM$\langle$P2]}}
                \texttt{explanation: the clue "gauge reading" refers to the measurement of a specific quantity, and "level" is a five{-}letter word that describes a measurement or a state of being at a certain point, fitting the description. Additionally, the feedback indicates that 'e' is in the second position, which matches the letter in "level".} \\
\texttt{agreement: yes} \\
            }
        }
    }
     \\ \\

    \theutterance \stepcounter{utterance}  
    & & \multicolumn{4}{p{0.6\linewidth}}{
        \cellcolor[rgb]{0.9,0.9,0.9}{
            \makecell[{{p{\linewidth}}}]{
                \texttt{\tiny{[P1$\langle$GM]}}
                \texttt{clue: gauge reading} \\
\texttt{guess\_agreement: yes} \\
\texttt{agreement\_explanation: the clue "gauge reading" refers to the measurement of a specific quantity, and "level" is a five{-}letter word that describes a measurement or a state of being at a certain point, fitting the description. Additionally, the feedback indicates that 'e' is in the second position, which matches the letter in "level".} \\
\\ 
\texttt{Provide your response only in this format.} \\
\texttt{explanation: details} \\
\texttt{guess: word} \\
            }
        }
    }
    & \\ \\

    \theutterance \stepcounter{utterance}  
    & \multicolumn{4}{p{0.6\linewidth}}{
        \cellcolor[rgb]{0.9,0.9,0.9}{
            \makecell[{{p{\linewidth}}}]{
                \texttt{\tiny{[P1$\rangle$GM]}}
                \texttt{explanation: The user confirmed that "level" is the correct answer. The feedback from the previous guess ("gauge") showed that 'e' is in the word but in the second position, which aligns perfectly with "level". Since the target word is confirmed as "level", no further guessing is required.} \\
\texttt{guess: level} \\
            }
        }
    }
    & & \\ \\

    \theutterance \stepcounter{utterance}  
    & & & \multicolumn{2}{p{0.3\linewidth}}{
        \cellcolor[rgb]{0.9,0.9,0.9}{
            \makecell[{{p{\linewidth}}}]{
                \texttt{\tiny{[GM$|$GM]}}
                \texttt{attempts = 2} \\
\texttt{target\_word = level} \\
\texttt{target\_word\_clue = gauge reading} \\
\texttt{guess = level} \\
\texttt{guess\_feedback = l<green> e<green> v<green> e<green> l<green>} \\
            }
        }
    }
    & & \\ \\

    \theutterance \stepcounter{utterance}  
    & & & \multicolumn{2}{p{0.3\linewidth}}{
        \cellcolor[rgb]{0.9,0.9,0.9}{
            \makecell[{{p{\linewidth}}}]{
                \texttt{\tiny{[GM$|$GM]}}
                \texttt{game\_result = WIN} \\
            }
        }
    }
    & & \\ \\

\end{supertabular}
}

\end{document}
